% CODEX: wstaw w 30b-geometrizacja-thurstona.tex pomiędzy końcem akapitu
% "Jej grupa podstawowa N_Z jest nilpotentna, lecz nieabelowa."
% a początkiem \textbf{Geometria $\mathrm{Sol}$.}.
% Nie powtarzaj definicji i tabeli krzywizn Nil z wcześniejszego akapitu.

\begin{przyklad}[Rozkład kontaktowy i holonomia w geometrii Nil]
\label{prz:thur-dop-nil-kontakt}
Forma $\theta=\dd z-x\,\dd y$ występująca w metryce Nil
wyznacza rozkład dwuwymiarowych podprzestrzeni
$\mathcal H=\ker\theta\subset TN$.
Ponieważ $\dd\theta=-\dd x\wedge\dd y$, mamy
\[
 \theta\wedge\dd\theta=-\dd x\wedge\dd y\wedge\dd z\ne0.
\]
Rozkład jest zatem \emph{kontaktowy}: nie jest styczny
do żadnej lokalnej rodziny powierzchni. Istotnie, gdyby
$\mathcal H$ było styczne do powierzchni, to nawiasy pól
stycznych również byłyby styczne. Tymczasem dla
$E_1=\partial_x$ oraz $E_2=\partial_y+x\partial_z$
mamy $E_1,E_2\in\mathcal H$ i
$[E_1,E_2]=\partial_z\notin\mathcal H$.

Na $N\to\R^2$, $(x,y,z)\mapsto(x,y)$,
forma $\theta$ jest także formą koneksji poziomej
przy pionowych translacjach współrzędnej $z$.
Poziome podniesienie krzywej $\gamma(s)=(x(s),y(s))$
spełnia $\theta(\dot{\widetilde\gamma})=0$, czyli
$\dot z=x\dot y$. Dla poziomego podniesienia zamkniętej
krzywej podstawy przyrost współrzędnej pionowej wynosi
\begin{equation}\label{eq:thur-dop-nil-holonomia-pola}
 \Delta z=\oint_\gamma x\,\dd y
 =\int_D\dd x\wedge\dd y,
\end{equation}
jeżeli $\gamma=\partial D$ z orientacją brzegową.
Ostatnia równość jest twierdzeniem Stokesa.
W szczególności obieg dodatnio zorientowanego
kwadratu jednostkowego zmienia $z$ o $1$.
Po przejściu do ilorazu przez kratę całkowitą jest to
okrągowe włóknienie nad $T^2$ o niezerowym skręceniu;
wartość bezwzględna jego liczby Eulera wynosi $1$.
Tak otrzymujemy geometryczny sens nieabelowego
komutatora w całkowitej grupie Heisenberga.
\end{przyklad}
